
\documentclass[12pt]{article}

% Layout.
\usepackage[top=1.2in, bottom=0.75in, left=1in, right=1in, headheight=1.0in, headsep=0pt]{geometry}

% Fonts.
\usepackage{mathptmx}
\usepackage[scaled=0.86]{helvet}
\renewcommand{\emph}[1]{\textsf{\textbf{#1}}}

% TiKZ.
\usepackage{tikz, pgfplots}
\usetikzlibrary{calc}
\pgfplotsset{my style/.append style={axis x line=middle, axis y line=middle, xlabel={$x$}, ylabel={$y$}}}
\pgfplotsset{compat=1.16}

% Misc packages.
\usepackage{amsmath,amssymb,latexsym,bm,array,multicol,enumitem}
\usepackage{graphicx}
\usepackage{xcolor}
\usepackage[most]{tcolorbox}
\usepgfplotslibrary{fillbetween}
\definecolor{solbg}{RGB}{225,240,252}
\definecolor{solframe}{RGB}{70,130,180}
\newtcolorbox{solution}[1][]{%
  breakable, enhanced,
  colback=solbg, colframe=solframe, boxrule=0.7pt, arc=3pt,
  left=6pt, right=6pt, top=5pt, bottom=5pt,
  before skip=6pt, after skip=6pt,
  fonttitle=\sffamily\bfseries\small, coltitle=white, title={Solution},
  #1}

% Commands to set various header/footer components.
\makeatletter
\def\doctitle#1{\gdef\@doctitle{#1}}
\doctitle{Use {\tt\textbackslash doctitle\{MY LABEL\}}.}
\def\docdate#1{\gdef\@docdate{#1}}
\docdate{Use {\tt\textbackslash docdate\{MY DATE\}}.}
\def\doccourse#1{\gdef\@doccourse{#1}}
\let\@doccourse\@empty
\def\docscoring#1{\gdef\@docscoring{#1}}
\let\@docscoring\@empty
\def\docversion#1{\gdef\@docversion{#1}}
\let\@docversion\@empty
\makeatother

% Headers and footers layout.
\makeatletter
\usepackage{fancyhdr}
\pagestyle{fancy}
\fancyhf{} % Clears all headers/footers.
\lhead{\emph{\@doctitle\hfill\@docdate} \medskip
\ifnum \value{page} > 1\relax\else\\
\emph{Name: \rule{3.5in}{1pt}\hfill \@docscoring}
\fi}

\rfoot{\emph{\@docversion}}
\lfoot{\emph{\@doccourse}}
\cfoot{\emph{\thepage}}
\renewcommand{\headrulewidth}{0pt}%
\makeatother

% Paragraph spacing
\parindent 0pt
\parskip 6pt plus 1pt

% A problem is a section-like command. Use \problem{5} for a problem worth 5 points.
\newcounter{probcount}
\newcounter{subprobcount}
\setcounter{probcount}{0}

\newcommand{\problem}[1]{%
\par
\addvspace{4pt}%
\setcounter{subprobcount}{0}%
\stepcounter{probcount}%
\makebox[0pt][r]{\emph{\arabic{probcount}.}\hskip1ex}\emph{[#1 points]}\hskip1ex}

\newcommand{\thesubproblem}{\emph{\alph{subprobcount}.}}

% Subproblems are an enumerate-like environment with a consistent
% numbering scheme.  Use \begin{subproblems}\item...\item...\end{subproblems}
\newenvironment{subproblems}{%
\begin{enumerate}[itemindent=8pt,leftmargin=0pt]%
\setcounter{enumi}{\value{subprobcount}}%
\renewcommand{\labelenumi}{\emph{\textsl{\alph{enumi})}} \,}%
}%
{%
\setcounter{subprobcount}{\value{enumi}}%
\end{enumerate}%
}

% Blanks for answers in normal and math mode.
\newcommand{\blank}[1]{\rule{#1}{0.75pt}}
\newcommand{\mblank}[1]{\underline{\hspace{#1}}}
\def\emptybox(#1,#2){\framebox{\parbox[c][#2]{#1}{\rule{0pt}{0pt}}}}

% Misc.
\renewcommand{\d}{\displaystyle}
\newcommand{\ds}{\displaystyle}
\newcommand{\threeopts}{{\small \hspace{-6mm} $\begin{matrix} \text{\textsc{converges}} \\ \text{\textsc{absolutely}} \end{matrix}$ \qquad\qquad $\begin{matrix} \text{\textsc{converges}} \\ \text{\textsc{conditionally}} \end{matrix}$ \qquad\qquad \textsc{diverges}} \bigskip}


\newcommand{\ba}{\mathbf{a}}
\newcommand{\bb}{\mathbf{b}}
\newcommand{\bc}{\mathbf{c}}
\newcommand{\bi}{\mathbf{i}}
\newcommand{\bj}{\mathbf{j}}
\newcommand{\bk}{\mathbf{k}}
\newcommand{\bu}{\mathbf{u}}
\newcommand{\bv}{\mathbf{v}}
\newcommand{\bw}{\mathbf{w}}

\newcommand{\ip}[2]{\mathrm{\left<#1,#2\right>}}
\newcommand{\vv}[2]{\mathrm{\left<#1,#2\right>}}
\newcommand{\vvv}[3]{\mathrm{\left<#1,#2,#3\right>}}


\doctitle{Math 252: Quiz 7 Solutions}
\docdate{29 October, 2026}
\doccourse{}
\docversion{Solutions}
\docscoring{\fbox{{\LARGE \strut}\hspace{0.8in} / 25}}

\begin{document}
30 minutes maximum.  No aids (book, calculator, etc.) are permitted.
Show all work and use proper notation for full credit.  Answers should
be in reasonably-simplified form.  25 points possible.

\begin{enumerate}

\item For each of the following series, determine whether the series
  converges or diverges. Answers without justification will receive no
  credit.

  \begin{enumerate}
  \item (5 points) $\ds \sum_{n=3}^{\infty} \frac{1}{2\sqrt{n}+5}$
    \begin{solution}
      Compare with $b_n = \dfrac{1}{\sqrt{n}}$. Both $a_n = \dfrac{1}{2\sqrt{n}+5}$ and $b_n$ are
      positive for $n \ge 3$, and
      \[
        \lim_{n\to\infty} \frac{a_n}{b_n} = \lim_{n\to\infty} \frac{\sqrt{n}}{2\sqrt{n}+5}
        = \lim_{n\to\infty} \frac{1}{2 + \frac{5}{\sqrt{n}}} = \frac12,
      \]
      which is finite and positive. $\ds\sum_{n=3}^\infty \frac{1}{\sqrt{n}}$ is a $p$-series with
      $p = \frac12 \le 1$, so it diverges. By the limit comparison test, the original series
      \boxed{\text{diverges}}.
    \end{solution}
    \vfill

  \item (5 points) $\ds \sum_{n=1}^{\infty} \sin^2(n) \left( \frac{1}{3} \right)^n$
    \begin{solution}
      Since $0 \le \sin^2(n) \le 1$,
      \[
        0 \;\le\; \sin^2(n) \left( \frac13 \right)^n \;\le\; \left( \frac13 \right)^n.
      \]
      $\ds\sum_{n=1}^\infty \left( \frac13 \right)^n$ is a geometric series with
      $|r| = \frac13 < 1$, so it converges. By the comparison test, the original series
      \boxed{\text{converges}}.
    \end{solution}
    \vfill

  \end{enumerate}

  \newpage
  
\item For each of the following series, determine whether the series
  converges absolutely, converges conditionally, or diverges. Answers
  without justification will receive no credit.

  \begin{enumerate}
  \item (5 points) $\ds \sum_{n=1}^{\infty} \frac{(-1)^n}{e^{1/n}}$
    \begin{solution}
      \[
        \lim_{n\to\infty} \left| \frac{(-1)^n}{e^{1/n}} \right| = \lim_{n\to\infty} \frac{1}{e^{1/n}}
        = \frac{1}{e^0} = 1,
      \]
      so the terms do not approach $0$ (they alternate near $\pm 1$). By the divergence test,
      the series \boxed{\text{diverges}}.
    \end{solution}
    \vfill

  \item (5 points) $\ds \sum_{n=3}^{\infty} \frac{(-1)^{n-1}}{2 \sqrt{n} +5}$
    \begin{solution}
      \emph{Absolute convergence.} $\ds\sum_{n=3}^\infty \left| \frac{(-1)^{n-1}}{2\sqrt{n}+5} \right|
      = \sum_{n=3}^\infty \frac{1}{2\sqrt{n}+5}$, which diverges by Problem 1(a).

      \emph{Alternating series test.} Let $b_n = \dfrac{1}{2\sqrt{n}+5}$. Then $b_n > 0$;
      $b_n$ is decreasing, since $2\sqrt{n}+5$ is increasing; and
      $\ds\lim_{n\to\infty} b_n = 0$. So the series converges.
      \[
        \boxed{\text{converges conditionally}}
      \]
    \end{solution}
    \vfill

  \item (5 points) $\ds \sum_{n=10}^{\infty} \frac{(-1)^{n+1}}{n^2 \ln(n)}$
    \begin{solution}
      For $n \ge 10$, $\ln(n) \ge \ln(10) > 1$, so
      \[
        0 \;<\; \left| \frac{(-1)^{n+1}}{n^2 \ln(n)} \right| = \frac{1}{n^2 \ln(n)} \;\le\; \frac{1}{n^2}.
      \]
      $\ds\sum_{n=10}^\infty \frac{1}{n^2}$ is a $p$-series with $p = 2 > 1$, so it converges.
      By the comparison test, $\ds\sum_{n=10}^\infty \frac{1}{n^2 \ln(n)}$ converges, so the
      series \boxed{\text{converges absolutely}}.
    \end{solution}
    \vfill
  \end{enumerate}

  \newpage

  BONUS (2 points) Does the series $\ds \sum_{n=1}^{\infty} \left( \frac{1}{n} -
    \cos\left(\frac{1}{n}\right) \right)$ converge or diverge? Answers
  without justification will receive no credit.
  \begin{solution}
    \[
      \lim_{n\to\infty} \left( \frac1n - \cos\left( \frac1n \right) \right) = 0 - \cos(0) = -1 \neq 0,
    \]
    so by the divergence test the series \boxed{\text{diverges}}.
  \end{solution}

\end{enumerate}

\end{document}

%-------------------------------------------------------------------------------------------------------------------------------------------------------------------------------------------------------------------

%%% Local Variables:
%%% mode: latex
%%% TeX-master: t
%%% End:
